\section{The Essence of the Traditions}

The New Age view of the various traditions run along the lines that ``all religions teach the same thing'', ``all paths are valid'', ``great teachers adapted religion for specific people'', and similar platitudes along the same lines of thought. From that, they seem to arrive at the illogical conclusion that if all paths lead to the summit, then there is no point in following a path. Even if we concede their assumptions, to get to the peak clearly requires following one of the traditional paths.

\textbf{Rene Guenon} himself makes this clear. Nevertheless, his fundamental assumption differs not much from the new age view, when he claims all traditions are the same; the same, too, for Schuon's
\emph{Transcendental Unity of Religions}. As anti-egalitarians, we cannot suddenly consider all traditions as somehow equal or equivalent. Rather, we need to grasp the inner essence of each tradition, i.e., what exact insight each one brings. Beyond platitudes, there is specificity, and it is that quality that separates the players from the patzers.

\paragraph{The Two Worlds}


From the Vedas to Plato to Augustine, we see the constant teaching of two worlds. Yet, to a great extent, those who claim a tie to Tradition as such seem to forget this teaching. This seems to me to be the case (readers can provide counterexamples) among neo-reactionaries and those who try to wed Tradition to various political movements. The focus is solely on the world of becoming, Even those---and these are legion relatively speaking---whose only exposure to Tradition is through \textbf{Julius Evola} cannot evade this. In the first paragraph of the first chapter of \emph{Revolt Against the Modern World}, he writes:

\begin{quotationx}
In order to understand both the spirit of Tradition and its antithesis, modern civilization, it is necessary to begin with the fundamental doctrine of the two natures. According to this doctrine there is a physical order of things and a metaphysical one; there is a mortal nature and an immortal one; there is the superior realm of ``being'' and the inferior realm of ``becoming.'' Generally speaking, there is a visible and tangible dimension and, prior to and beyond it, an invisible and intangible dimension that is the support, the source, and true life of the former.

Anywhere in the world of Tradition, both East and West and in one form or another, this knowledge (not just a mere theory) has always been present as an unshakeable axis around which everything revolved. Let me emphasize the fact that it was knowledge and not theory. As difficult as it may be for our contemporaries to understand this, we must start from the idea that the man of Tradition was aware of the existence of a dimension of being much wider than what our contemporaries experience and call ``reality''.
\end{quotationx}


Hence, the man of Tradition must strive to gain true knowledge of that realm of being, even to the point of obsession. Of course, that realm is not obvious, whence the need for Tradition and true spiritual teachings. The natural man is aware only of the inferior world of becoming. That world appears to be governed by evolution and the will to power, and man's only purpose is the pursuit and fulfilment of sensory pleasures. We don't deny those appearances, rather we deny their ultimacy as an explanation of the world process.

To get back to the point, we can find the essence of the various traditions in how they understand and teach the higher world. \textbf{Valentin Tomberg}, in the \emph{Covenant of the Heart}, scatters interesting insights about the role of the various traditions in salvation history. I will draw on some of them. First, this is Tomberg's definition of salvation:

\begin{quotex}
Salvation is the helping hand offered to idealistically minded human beings who, though no longer identifying themselves with natural evolution, still lack sufficient strength of themselves to alter its course.
\end{quotex}


To rephrase it in Guenon's terms, the actualization of man's possibilities as salvation is not possible within the material world process, but must arise from a transcendent source. So there is a salvation history as that ``helping hand'' revealed itself in various ways. Tomberg writes:

\begin{quotex}
The hand of God stretched down form on high to human beings mired in evolution's stream so that they may be raised up and helped to return to the divine archetypal world ... which came in the form of successive stages of revelation, descending like tumbling cascades of a waterfall, awakening, the consciousness of mankind to its divine archetype. This process of revelation took place through seers and prophets and was subsequently recorded and preserved in holy scriptures. Thus, for instance, the revelation vouchsafed to the seven Holy Rishis forms the basis of the Vedas, the Holy Writ of India. The Zend-Avest of ancient Iran records the revelation communicated to Zarathustra when, on a high mountain fastness, he confronted the God of Light---Ahura Mazda---face to face. The Bible tells of the revelation received by Moses on Mount Sinai, where he encountered, also face to face, the God who bore the ineffable name of YHWH.
\end{quotex}


Note that these are all revelations of the superior world of Being. Regarding Brahmanism, Tomberg writes:

\begin{quotex}
The primordial revelation of India announces the good news that the divine-archetypal Man, the true self of mankind, is not the empirical self, and likewise that the world---the real world that is, the divine-archetypal world---is not the empirical world of natural evolution.
\end{quotex}


I was speaking to two Hindus today about spiritual matters and, to my surprise, they could not see this as good news. India is not as spiritual as Guenon believed, but it full of charlatans charging money for quack cures and instant enlightenment; it should be no surprise that many Westerners are drawn to such nonsense. For example, one such guru says your karma will be burned off by walking X times around a certain temple which reciting a mantra he gives you. I told them that salvation cannot be achieved by such mechanical means. Unfortunately, the doctrines of karma and reincarnation seem to be traps and the opposite of good news. The two were obsessed about how to burn up karma and achieve a better life, or actually a better next life.

It reminded me more of the medieval buying of indulgences than any free movement of the spirit. Their idea is that a good life would be free of the many trials we undergo. Yes, I would agree if understood properly. The trials are the results of man's fall and are part of the world of becoming, not of being. Christian salvation promises to restore man to the state before the fall, even in an instant.

One of them told me he was depressed at the thought of the number of future lives that are necessary for him to achieve deliverance. That surprised me because the spiritual life should be one of joy. If one is still obsessed only with improving one's material circumstances in the world of becoming, then, yes, depression will result unless one sinks into in full and unconsciously.

I told him that one's true self transcends the empirical self and that the solution is vertical, not horizontal. As one learns to detach and observe life, emotions, and thoughts dispassionately, there will be increasingly long states of bliss.

To move on, Tomberg then writes:

\begin{quotex}
The good news of Zarathustra was that the world and Man represented an admixture of two distinct world orders---that of the principle of light and that of the principle of darkness, or the divine--archetypal world order and that of natural evolution---and that the latter would in the end be vanquished by Soshyans ``who through will overcomes death'', to be followed by the resurrection of the dead.
\end{quotex}


Note how different this is from ``academic'' and ``learned'' explanations of Zoroastrianism. They have no clue that Zarathustra was talking about the worlds of being and becoming. Moving on to the next revelation, Tomberg continues:

\begin{quotex}
Finally, the good news revealed to Moses and the prophets proclaims that mankind, having fallen into the serpent's domain through the Fall, can yet be redeemed, for the very essence and epitome of the divine-archetypal world order Himself was to become Man so that mankind might thereby triumph over the domain of the serpent and the dead might resurrect.
\end{quotex}


This has set the stage, finally, for our ultimate purpose, to wit, the self-understanding of the medieval Tradition, how it relates to other traditions, the source of its metaphysical teachings, and how it developed organically from the Roman religion that preceded it. God willing, this will appear this weekend.


\flrightit{Posted on 2014-02-26 by Cologero}

\begin{center}* * *\end{center}

\begin{footnotesize}
\begin{sffamily}

\texttt{August on 2014-02-28 at 02:46 said:}

I wouldn't call recognition of the convergence of paths in transcendental unity an ``egalitarian'' or levelling view, so long as we understand transcendence in the absolute sense.

Guenon, for example, did differentiate between traditions, not so much horizontally, but vertically, which allows comparison not only of traditions, but of their inner accommodations and followers. To compare them horizontally is not the most useful exercise; the chief reason different paths exist is because races and ages are qualitatively distinct, so that the variety is necessary, and does not need to be reduced by humans. And, like any other contingency, traditions are capable of taking an indefinite number of forms, so that studying the insights of those before us can only ever, at most, bring one to the threshold of the Absolute, the purpose of every tradition being fulfilled at that point.

Consider Guenon here, discriminating:

`` The preceding will suffice to mark the profound distinction that exists between the Great Architect of Masonry on the one hand and on the other hand the gods of the various religions, who are only so many aspects of the Demiurge. Moreover, it is incorrect to identify... the anthropomorphic God of exoteric Christians with Jehovah... or with Allah... ''

\hfill

\texttt{August on 2014-02-28 at 03:14 said:}

As exit says, the new age thought will tend to reject all valid paths, because it thinks that is has found, or is about to find, something better, a `true key', or some great synthesis that operates to combine the good things and skim off the bad all at once. We see this for what it is, yet another combination borne of desperation, stirred by a whiff of the esoteric that whets the collective appetite of a multitude that cannot seize the jewel as a multitude.

But then there are others, certainly not new age, who see (exoteric) tradition, especially religion, as a mainly social phenomenon, and who have a good grasp of transcendence, even if it only be cognitive, enough to create perspective. Combined with a certain attitude fashioned in the exasperation of today, these types will assert a freedom that can participate in any tradition or none, to whatever extent they can or wish. Why shouldn't they? Sure, this world is wrong, but they are its products; and there is a middle path between succumbing to error and adopting a regular lifestyle today, and the strained reversion fantasy that preoccupies many a `traditionalist'. That middle path, for some, is treating religion as a burden to be avoided without sufficient incentive to participate, just like anything else in our corner of the kingdom of becoming. And like any burden, setting it aside is not all gain, but one thinks about this before deciding.

\hfill

\texttt{Cologero on 2014-02-28 at 08:51 said:}

August wrote: ``Guenon, for example, did differentiate between traditions, not so much horizontally, but vertically, which allows comparison not only of traditions, but of their inner accommodations and followers.''

Precisely, August. But it is one thing to assert its possibility, and quite another to actually do the comparison. That is what we are attempting. That Guenon quote is rather vacuous; which Christian Church Father proclaimed an ``anthropomorphic God''. I am getting weary of broad strokes ... where have we indulged in a ``strained reversion fantasy'', August? Do you just like the sound of your ``voice'' or do you actually want to engage in a conversation?

\hfill

\texttt{August on 2014-02-28 at 20:40 said:}

I didn't have Gornahoor in mind when I mentioned reversion fantasies. This article was strong, and I look forward to the following.

Don't pay me too much mind, I don't know what I want. Carry on, and if the time comes to banish me, say the world and I'll say no more.

\hfill

\texttt{JA on 2014-03-01 at 10:33 said:}

Personally when I read a book I refrain from making final judgement until I've reached the end; as such I will wait until post 1001 before I can say whether I fully agree with Gornahoor or not, until then all I will do is mention that so far it's been extremely helpful and enlightening to me.

\hfill

\texttt{JA on 2014-03-01 at 10:37 said:}

New age teaching as I understand claims that the forms (exotericism) is meaningless and the core (esotericism) is the same everywhere; the same approach Evola has -- and also Crowley. Dont forget that before encountering Guenon's ideas Evola was involved in the Italian new age milieu albeit those of higher intelligence than the simple frauds that trick old ladies.

Guenon even if he never quite stated it this way saw Hinduism as the highest form of Tradition, but since it's a closed racial path, he saw Islam as the next best thing. What Cologero is showing us with Tomberg is a different way of looking at things which Guenon never considered and I am finding myself in agreement with him.

\hfill

\texttt{Scardanelli on 2014-03-01 at 12:03 said:}

Do not despair brother August. The ego must be checked from time to time so that we may experience what is beyond the chattering of the mind. For what it's worth, I don't think it's a simple matter of agreeing or disagreeing with the content here, as that is not enough and would be rather presumptuous besides. Truth is not understood but seen, and Gornahoor has outlined the means to apprehend this truth within. Until this becomes active within us, we will not know what to do with ideas; we will not know how to ``handle'' them. Otherwise we can continue to comment and banter and discuss, or contribute little one liners, and we will remain destitute men of this world without knowledge of the kingdom of God.

``Happy is the man whom truth by itself doth teach, not by figures and transient words, but as it is in itself.'' -Kempis

I'm glad that you are back in good health Cologero. Your wisdom has been sorely missed. I will pray for your continued health (though perhaps somewhat selfishly).


\hfill

\end{sffamily}
\end{footnotesize}

